\begin{lemma}
\label{lemma-H1-regular-quasi-regular}
An $M$-$H_1$-regular sequence is $M$-quasi-regular.
\end{lemma}

\begin{proof}
Let $R$ be a ring and let $M$ be an $R$-module.
Let $f_1, \ldots, f_r$ be an $M$-$H_1$-regular sequence.
Denote $J = (f_1, \ldots, f_r)$. The assumption means that we have
an exact sequence
$$
\wedge^2(R^r) \otimes M \to R^{\oplus r} \otimes M \to JM \to 0
$$
where the first arrow is given by
$e_i \wedge e_j \otimes m \mapsto (f_ie_j - f_je_i) \otimes m$.
Tensoring the sequence with $R/J$ we see that
$$
JM/J^2M = (R/J)^{\oplus r} \otimes_R M = (M/JM)^{\oplus r}
$$
is a finite free module. To finish the proof we have to prove
for every $n \geq 2$ the following: if
$$
\xi = \sum\nolimits_{|I| = n, I = (i_1, \ldots, i_r)}
m_I f_1^{i_1} \ldots f_r^{i_r} \in J^{n + 1}M
$$
then $m_I \in JM$ for all $I$. In the next paragraph, we prove
$m_I \in JM$ for $I = (0, \ldots, 0, n)$ and in the last paragraph
we deduce the general case from this special case.

\medskip\noindent
Let $I = (0, \ldots, 0, n)$. Let $\xi$ be as above. We can write
$\xi = m_1 f_1 + \ldots + m_{r - 1}f_{r - 1} + m_I f_r^n$.
As we have assumed $\xi \in J^{n + 1}M$, we can also write
$\xi = \sum_{1 \leq i \leq j \leq r - 1} m_{ij}f_if_j +
\sum_{1 \leq i \leq r - 1}m'_i f_if_r^n + m'' f_r^{n + 1}$.
Then we see that
$$
\begin{matrix}
(m_1 - m_{11}f_1 - m'_1f_r^n)f_1 + \\
(m_2 - m_{12}f_1 - m_{22}f_2 - m'_2f_r^n)f_2 + \\
\ldots + \\
(m_{r - 1} - m_{1 r - 1}f_1 - \ldots - m_{r - 1 r - 1}f_{r - 1}
- m'_{r - 1}f_r^n)f_{r - 1} + \\
(m_I - m'' f_r)f_r^n = 0
\end{matrix}
$$
Since $f_1, \ldots, f_{r - 1}, f_r^n$ is $M$-$H_1$-regular by
Lemma \ref{lemma-mult-koszul-regular} we see that $m_I - m'' f_r$
is in the submodule $f_1M + \ldots + f_{r - 1}M + f_r^nM$.
Thus $m_I \in f_1M + \ldots + f_rM$.

\medskip\noindent
Let $S = R[x_1, x_2, \ldots, x_r, 1/x_r]$. The ring map $R \to S$ is faithfully
flat, hence $f_1, \ldots, f_r$ is an $M$-$H_1$-regular sequence in $S$, see
Lemma \ref{lemma-koszul-regular-flat-base-change}.
By
Lemma \ref{lemma-change-basis}
we see that
$$
g_1 = f_1 - \frac{x_1}{x_r} f_r,
\ \ldots,
\ g_{r - 1} = f_{r - 1} - \frac{x_{r - 1}}{x_r} f_r,
\ g_r = \frac{1}{x_r}f_r
$$
is an $M$-$H_1$-regular sequence in $S$. Finally, note that our element
$\xi$ can be rewritten
$$
\xi = \sum\nolimits_{|I| = n, I = (i_1, \ldots, i_r)}
m_I (g_1 + x_1 g_r)^{i_1} \ldots (g_{r - 1} + x_{r - 1} g_r)^{i_{r - 1}}
(x_rg_r)^{i_r}
$$
and the coefficient of $g_r^n$ in this expression is
$$
\sum m_I x_1^{i_1} \ldots x_r^{i_r}
$$
By the case discussed in the previous paragraph this sum is in
$J(M \otimes_R S)$. Since the monomials $x_1^{i_1} \ldots x_r^{i_r}$
form part of an $R$-basis of $S$ over $R$ we conclude that $m_I \in J$
for all $I$ as desired.
\end{proof}